\documentclass[10pt,a4paper]{article}
\usepackage{onepager}

\begin{document}
\small

\ophead{The Process: A First-Class Analytical Unit}
{Conditional-Information Signal Selection and the Effective Degrees of Freedom of an Alpha Book}
{\code{process_preprint.tex} $\cdot$ 14pp $\cdot$ Jul 2026 $\cdot$ renamed from \code{process_related}
 2026-08-08 $\cdot$ methods/framework paper, empirics stated as future work}

\begin{multicols}{2}

\opsec{Thesis}
A microstructure feed produces hundreds of features, most mutually redundant. The operative question is
therefore \textbf{not} whether feature $X$ predicts, but \textbf{how much non-redundant information the
set carries in aggregate}, and \textbf{how many statistically independent bets that supports}.

\opsec{The unit}
Beside the \textbf{feature} (what is computed) and the \textbf{algorithm} (how it trades), place a third
first-class unit: the \textbf{process} — a typed, registry-bound analytical contract that certifies
whether a feature, or a derivative/combination of features, carries information about a
\emph{precisely specified} target. Every process ships behind a mandatory planted-signal conformance gate.

\opsec{Why information, not correlation}
Mutual information as the currency of prediction, justified by \textbf{Fano's inequality} — information is
a hard \emph{lower bound} on achievable error. A feature with zero MI about the target cannot be traded by
any model, however clever; correlation gives no such guarantee.

\opsec{The central quantity}
Fundamental law of active management:
\[
  \boxed{\ \IR \approx \ICbar\sqrt{\Neff}\ },\qquad
  \Neff=\frac{N}{1+(N-1)\rhobar}
\]
\textbf{The binding quantity is the stress-regime correlation, not the calm-regime one.} A book of 20
signals with $\rhobar=0.1$ in calm and $0.8$ in stress has $\Neff\approx6.9$ when nothing is happening and
$\Neff\approx1.2$ exactly when it matters. Reporting calm-regime breadth is the standard way to lie to
yourself about diversification.

\opsec{Five orthogonalization mechanisms, each a process}
\begin{opnums}
\item \textbf{Conditional-MI greedy selection} — the core extractor. Add the feature maximising
      $\MI(X;Y\mid S)$ given what is already selected.
\item \textbf{Interaction-information synergy mining} — finds pairs that carry jointly what neither
      carries alone (negative interaction information); the thing greedy marginal screens structurally
      cannot see.
\item \textbf{Training-prefix residualization} — pure-innovation extraction, fit on the prefix only so
      the residual is causal.
\item \textbf{Marchenko--Pastur-denoised PCs} — separates real factors from the eigenvalue bulk that
      pure noise produces at finite $T/N$. Tells you how many factors are \emph{real}.
\item \textbf{Capstone: a decorrelated signal book that reports its own honest $\Neff$.}
\end{opnums}

\opsec{Making ``price action'' explicit}
An aspect $\times$ horizon target taxonomy — a process must name \emph{which} aspect (sign, magnitude,
volatility, timing) at \emph{which} horizon it certifies. Half the ambiguity in feature research is a
target that was never written down.

\opsec{Minting new degrees of freedom}
A derivative-operator layer creates new feature DoF from existing ones: fractional differentiation
(memory-preserving stationarity), leaky integration, causal spectral features — all under a
\textbf{program-level FDR ledger}, because minting features is exactly how a search silently multiplies
its own multiple-testing burden.

\opsec{Estimation hygiene}
Keeping the breadth estimate honest is its own section: $\Neff$ computed from a noisy $\hat\rho$ is
biased upward, so the paper treats the estimate itself as something requiring a gate rather than a
formula you evaluate once and quote.

\opsec{Implementation status (NAT)}
PROC critical path complete 2026-07-30: null-calibration, \code{conditional_predictability}, standing
3-bar eval, process-FDR + cross-run ledger, \code{horizon_label_scan}, predictability surface
(\code{nat viz predictability}).
\caution{First real surface: 284 BTC cells, argmax $z=13.2$ but $q=0.56$ $\Rightarrow$ correctly
\textsf{[SPEC]}, \textbf{0 discoveries}. BH at the 100-shuffle $p$-floor over 284 cells \emph{cannot}
pass — real discoveries need deeper nulls ($n_{\text{shuffles}}\gtrsim m/\alpha$) on focused candidate
sets, or more clean data. The zero is a power problem, not a verdict on the features.}

\opsec{The two kinds (they share registry, CLI, persistence)}
\textbf{Evaluation process} — scores existing columns, produces no new series:
\[
  \mathcal{E}:(\mathbf{D},\mathcal{C})\longmapsto\mathcal{R},\qquad \mathbf{D}\in\R^{T\times p}
\]
\textbf{Transform process} — mints $m$ new series on the input index (PCs, residuals, fractional
differences):
\[
  \mathcal{T}:(\mathbf{D},\mathcal{C})\longmapsto(\mathbf{D}',\mathcal{R}),\qquad \mathbf{D}'\in\R^{T\times m}
\]
Context $\mathcal{C}$ carries symbol, timeframe, price column, horizon set, cost parameters, target, and
auxiliary sources. Identified by unique name, registered by decorator, parameterized by a config section
keyed to that name, persisted through \emph{one} result schema. The uniformity is the point: a new
analysis is a registry entry, not a new script.

\opsec{The planted-signal discipline}
Every unit ships behind a mandatory planted (synthetic) conformance gate \textbf{written before the
implementation}. Not a style preference — three estimator bugs in Stage 1 were caught only this way, and
nothing else would have caught them, because on real data a subtly wrong estimator still returns
plausible numbers.

\opsec{Why this is a paper and not just infrastructure}
It supplies the theory that says what the other papers' numbers \emph{mean}: \#2's five gates are a
special case of a process contract, \#7's eight independent axes are an $\Neff$ claim, \#5's whole
argument for an exogenous signal is an $\Neff$ argument, and \#8's gate ladder is this discipline applied
one level down, to data rather than signals.

\end{multicols}

\opfoot{One-pager for personal use. The sentence to remember: \emph{the binding correlation is the
stress-regime correlation}. Everything else here is machinery for measuring it honestly.
Detail: \code{docs/specs/process_layer.md} (PROC-1..18).}

\end{document}
